\section{Discussion}
\label{sec:discussion}

Guardian Lens shows that a compact black-box auditor can distinguish
researcher-controlled visual-conditional and generalized decision
policies from their observable behavior. The primary evidence is not
classification accuracy alone. It is the convergence of matched visual
counterfactuals, explicit costly decisions, profile-specific behavioral
signatures, a calibration-only classifier, and predictions frozen before
profile unblinding. Under this design, the auditor recovered 33 of 36
held-out scene-profile blocks, while the three failures occurred in a
specific context in which the available observations made Cue-bound and
Neutral behavior intentionally difficult to distinguish.

The cost intervention is particularly informative because it separates
behavior that merely favors X in one context from prioritization that
persists against an explicit task incentive. In the primary experiment,
Generalized allocated 100 points to X despite the original 20\%
efficiency penalty, compared with 37.50 for Cue-bound and 1.24 for
Neutral. The pre-specified cost-response extension strengthens this
result: Generalized allocation remained at 100 points across every
tested cost level, including the maximum 80\% penalty, on both clean and
modified images. Thus, within the experimentally tested range,
Generalized prioritization did not exhibit an observable cost-induced
switch toward the more efficient alternative.

Cue-bound behavior provides the complementary result. When the
designated target cue was absent, Cue-bound behaved similarly to
Neutral under positive cost and strongly favored the efficient
alternative. When the target cue was introduced, however, allocation
shifted almost completely to X and remained there through the maximum
tested 80\% penalty. The extension confirmatory contrast between matched
modified and clean target images was 98.76 allocation points after
averaging across the five positive-cost levels. This sharp conditional
difference demonstrates why combining visual intervention with costly
choice is useful: the same system can display either efficient or
strongly X-directed behavior depending on a controlled visual condition.

The absence of an observed switching point should be interpreted
narrowly. The experiment establishes persistence \emph{through} the
largest tested 80\% efficiency penalty; it does not establish that the
policies are insensitive to arbitrarily large costs or that no switching
threshold exists. Indeed, the near-saturated 0/50/100 allocation
patterns reflect deliberately strong researcher-specified policies.
Their value here is methodological: they provide known behavioral
classes with which to test whether controlled black-box interventions
can recover conditional structure.

The prompt-surface extension addresses a complementary internal-validity
concern. The primary and cost-response experiments alone leave open
whether the observed policy separation depends materially on incidental
properties of a particular prompt realization. Reversing organization
presentation order, exchanging the superficial X/Y label-role assignment,
reversing JSON output-key order, and using semantically equivalent task
wording did not eliminate either pre-specified confirmatory signature.
All ten variant-specific Generalized-versus-Neutral and Cue-bound
modified-versus-clean contrasts remained positive and significant after
joint Holm correction.

The result should not be interpreted as perfect numerical invariance.
The largest deviation from contemporaneous canonical P0 was 14.81
allocation points for Cue-bound target activation under reversed JSON-key
order. Thus, response formatting can affect effect magnitude even when
the qualitative behavioral distinction remains intact. Within the four
tested transformations, however, no prompt-surface manipulation erased
or reversed the policy signatures. This strengthens the interpretation
that the observed distinctions reflect more than one particular ordering,
label assignment, output-key order, or wording realization.


The distractor condition exposes the complementary limitation of
passive observation. Cue-bound does not systematically activate for the
visually similar distractor, supporting specificity to the designated
target within this benchmark. Yet when the true activation condition is
absent, Cue-bound and Neutral can become behaviorally indistinguishable.
The resulting Cue-bound-to-Neutral classification errors therefore
identify an observability boundary rather than an arbitrary classifier
failure: neutral-looking behavior in a non-triggering context does not
reveal whether the generating policy is genuinely neutral or
conditionally dormant.

This motivates active black-box auditing. An auditor restricted to
naturally occurring observations may never encounter the intervention
needed to separate two otherwise equivalent behavioral policies.
Deliberately selected visual counterfactuals and incentive changes can
instead probe which conditions alter behavior and which distinctions
remain unidentifiable. Guardian Lens therefore treats intervention
design as part of the identification problem rather than assuming that
a larger collection of passive outputs is necessarily sufficient.

Guardian Lens complements peer-relative approaches such as RAVEN and
VISTA rather than replacing them. Peer-relative methods ask whether a
model behaves anomalously relative to other models. Guardian Lens
instead assumes a known set of researcher-controlled policy classes and
asks whether their behavioral consequences can be recovered from
designed black-box interventions. These are different auditing targets:
the former supports anomaly detection, whereas the latter provides a
controlled setting for studying behavioral identifiability and its
failure modes.

More broadly, all results should be interpreted operationally. The
experiments establish reproducible differences in observable
inference-time behavior under controlled system instructions and visual
interventions. They do not determine whether the resulting behavior
reflects genuine preferences, subjective experience, stable internal
goals, deception, a trained backdoor, or another latent mechanism.
